\section{训练协议与重要修订}\label{app:protocols}

\subsection{语料与匹配学生模型}
PTB 和 WikiText-2 使用归档的原始划分；TinyStories 使用种子 42 将训练数据按 98/2 划分，并将原始验证集用作测试集。PTB 包含 42,068 个非空训练文本行和 4,439 个训练数据块。TinyStories 使用 300,000 个训练文本行（截断前 67,081,279 个词元；262,036 个数据块）、42,391 个验证文本行（9,440,302 个词元；36,876 个数据块）和 21,990 个测试文本行（4,765,822 个词元；18,616 个数据块）。固定长度为 256 个词元的数据块在评估时使用位置 1--255；留出集的目标词元不进入训练。

每个主要实验种子独立训练自己的 $S_0$，使用 20 轮训练、批大小 32、学习率 $2\times10^{-4}$、权重衰减 0.01、AMP，并按验证集 NLL 选择检查点。初始训练使用 AdamW、源实现的默认配置和固定学习率。续训重新初始化 AdamW，使用 10 轮训练、批大小 32、学习率 $10^{-4}$、权重衰减 0.01、5\% 线性预热后接余弦衰减、AMP、裁剪阈值 1 和 $\tau=2$。CE（0）与 KD（5）之间仅 $\lambda$ 不同。前向 KL 对有效目标词元取平均，并包含温度平方缩放。这些配置在本轮论文修改前已经冻结。

共同的 WikiText-2 子集包含 512 个验证数据块，使用种子 20260920 无放回抽样后排序（130,560 个目标词元）。完整验证集上的 TT/TG 指标保持单独报告。J 在该子集上的重复评估值 4.300644 和 4.300646 保留原始精度，不合并取平均。

\subsection{教师模型训练与准入条件}
PTB 和 TinyStories 的教师模型均将宽度为 256 的六层 Transformer 与宽度为 240、降维宽度为 64、时间间隔为 $\{1,2,4\}$ 的六层 Grassmann 分支组合；两个分支均使用 dropout 0.1。源模型训练使用种子 42、20 轮训练、批大小 32、序列长度 256 和学习率 $3\times10^{-4}$；四个源模型均按验证集选择第 20 轮。融合教师模型联合训练十轮，在第 1 轮冻结分支，使用分支学习率 $10^{-5}$、权重衰减 0.01、AMP、初始 $\alpha=0.5$ 和 late-$k=1$。两者的融合权重学习率分别为 $5\times10^{-3}$ 和 $10^{-2}$。验证集分别选择融合模型的第 8 和第 10 轮；主要续训保留学得的 $\alpha=0.470293$ 和 $0.599406$，而强制 $\alpha=0.5$ 适用于 WikiText-2 对照。

表~\ref{tab:teacher_training_protocol} 记录了教师模型的训练模式。J 和 TT 联合训练十轮；A 仅训练融合权重五轮。状态干预固定部署时的架构与组成，而不固定教师模型的训练历史。TT 的源 Transformer 具有六层网络、宽度 256、dropout 0.1、批大小 32、学习率 $3\times10^{-4}$ 和 20 轮训练，使用种子 42 和 123；两者均选择第 12 轮。TT 联合训练选择第 10 轮。其融合学习率与历史 TG 配置不同。
\begin{table*}[!t]\centering
\begin{tabular}{llrrr}
\toprule
教师模型 & 模式 & 训练轮次预算 & 分支学习率 & 融合权重学习率\\
\midrule
J/TG & 联合训练 & 10 & $10^{-5}$ & $5\times10^{-3}$\\
A & 仅训练融合权重 & 5 & 冻结 & $10^{-2}$\\
TT & 联合训练 & 10 & $10^{-5}$ & $10^{-2}$\\
\bottomrule
\end{tabular}
\caption{归档的 WikiText-2 教师模型训练配置。所有模型均以 $\alpha=0.5$ 初始化，使用批大小 32、权重衰减 0.01 和 AMP，并按验证集选择检查点。联合训练分支在第 1 轮冻结；A 全程冻结分支。评估和 KD 均强制 $\alpha=0.5$。这些训练历史不能单独识别教师模型的标量质量或架构效应。}
\label{tab:teacher_training_protocol}
\end{table*}


联合训练改变了 179 个 J/A 状态张量，全部状态张量共计 191 个。缩减数据的准入 smoke test 中，裁剪比例为 1，AMP/非有限值比例为 0.16。经授权，以全量数据进行的一轮训练准入检查替代了该测试，目标函数与裁剪阈值均未改变。J/A 准入检查的裁剪比例为 0.234694 和 0.275510，两者的 AMP/非有限值比例均为 0.013605。A/J 参数梯度比 1.102792、1.148522 和 1.093056 均满足预注册区间 $[0.5,2]$。A 正式训练第一轮的裁剪比例为 0.666667、0.503401 和 0.513605；AMP/非有限值比例的最大值为 0.013605。

T-only 需要单独授权的准入条件修订：全量数据准入检查中每次更新均触发裁剪，但参数梯度保持有限，T/F 比为 1.3164、1.2788 和 1.2324，AMP/非有限值比例为 0.013605。正式训练保留了全部配置。第一轮仍完全裁剪饱和；各轮次的平均裁剪比例为 0.4990、0.4605 和 0.4398，在训练后期降至约 0.0034。三个种子均选择第 10 轮。永久保留的 \texttt{CLIP\_SATURATION\_WARNING} 限制了对 F--T 差距的解释。

TT 学生模型第一轮的裁剪比例为 0.928571、0.843537 和 0.874150；对应的跨训练轮次均值为 0.170068、0.137755 和 0.127551。AMP/非有限值比例的最大值为 0.013605。这些轨迹与 TG 并非严格的优化匹配。

\subsection{TinyStories 稳定性与接受的重试}
主要 TinyStories 实验的所有 CE 和 KD 分支均选择第 10 轮。各种子的 CE 平均裁剪比例均为 0.000391；KD 均值为 0.236891、0.283344 和 0.240090（表~\ref{tab:tinystories_optimization}）。AMP/非有限值比例的最大值在 CE 中为 0.000488，在 KD 中为 0.000611。接受的终点指标均有限，KD 并未持续完全裁剪饱和。

种子 456 的一次运行因基础设施功率限制在第 2 轮中断，并在产生任何最终终点前被排除。经授权完成的重试从相同 $S_0$ 重新开始，使用新优化器，并保持教师模型、种子、数据、目标函数、预算和选择规则不变。只有完成的重试计入结果。实验使用 RTX 3090 GPU；详细的总体计算资源统计不可用，本文不对此进行推断。
\begin{table}[!t]\centering
\begin{tabular}{lrrrrr}
\toprule
种子 & 残余优势 & KD KL & CE 范数 & KD 范数 & KD 裁剪比例 \\
\midrule
42 & -0.021960 & 0.380224 & 0.486891 & 0.973546 & 0.236891 \\
123 & -0.021422 & 0.381402 & 0.488263 & 0.981438 & 0.283344 \\
456 & -0.019714 & 0.383282 & 0.483211 & 0.973942 & 0.240090 \\
\bottomrule
\end{tabular}
\caption{主要 TinyStories 实验的训练监测指标。残余优势为教师模型相对于 CE 测试集终点的优势。KL 包含温度平方缩放；范数和裁剪比例为已完成运行各训练轮次的均值。}\label{tab:tinystories_optimization}
\end{table}

